\documentclass[a4paper, 11pt]{article}
%\documentclass{article}
\usepackage[utf8]{inputenc}
\usepackage{hyperref}
\usepackage{calc}
%\usepackage{fourier} %% changes everything to sans-serif!
\usepackage{amssymb}
\usepackage{amsmath}
\usepackage{amsthm}
\usepackage{amsfonts}
\usepackage{mathtools}
\usepackage{textcomp}
%\usepackage{tikz}
\usepackage{tikz-cd}
\usetikzlibrary{cd}
\usepackage{dirtytalk}
\usepackage{titlesec}
\usepackage{verbatim}
\usepackage{pdfpages}
\usepackage{caption}
%%\captionsetup[figure]{labelfont=footnotesize}
\usepackage{hyperref}
\usepackage{parskip}

%% neu
\usepackage{physics}
\usepackage{accents} %% for multiple subscripts
\usepackage{fancyhdr} %% for header and footer
\usepackage[a4paper, left=2.5cm, right=2.5cm, top=2cm]{geometry}
\usepackage{graphicx} % Required for inserting images
%\usepackage[ngerman]{babel}
\usepackage{csquotes}
\MakeOuterQuote{"}

\title{Tuning Systems Exercises}
\author{Catharina Hock, Stefanie Hövermann}
\date{August  2026}

\newcommand{\answerline}[1][0.5cm]{%
    \noindent\rule{#1}{0.4pt}
}

\begin{document}
\maketitle 

\section*{Exponential Functions and Logarithms}

\subsection*{Exercise 1: Powers}

Calculate the following.

\begin{enumerate}
    \item $2^5=$ \rule{2cm}{0.4pt}
    \item $10^3=$ \rule{2cm}{0.4pt}
    \item $5^0=$ \rule{2cm}{0.4pt}
    \item $3^4=$ \rule{2cm}{0.4pt}
    \item $2^{-2}=$ \rule{2cm}{0.4pt}
\end{enumerate}

\vspace{0.5cm}

\subsection*{Exercise 2: Exponential Growth}

A bacteria culture starts with 200 bacteria. Every hour the number doubles.

\begin{enumerate}
    \item Complete the table.

    \begin{center}
    \begin{tabular}{c|ccccc}
        Time (h) & 0 & 1 & 2 & 3 & 4\\
        \hline
        Bacteria & 200 & \rule{1cm}{0.4pt} &
        \rule{1cm}{0.4pt} &
        \rule{1cm}{0.4pt} &
        \rule{1cm}{0.4pt}
    \end{tabular}
    \end{center}

    \item Write a formula for the number of bacteria after $t$ hours.

    \vspace{1cm}

    \item How many bacteria are there after 6 hours?

    \vspace{1cm}
\end{enumerate}

\subsection*{Exercise 3: Exponential Decay}

A radioactive substance loses half of its mass every day. Initially it has 80 g.

\begin{enumerate}
    \item Complete the table.

    \begin{center}
    \begin{tabular}{c|ccccc}
        Days & 0 & 1 & 2 & 3 & 4\\
        \hline
        Mass (g) & 80 &
        \rule{1cm}{0.4pt} &
        \rule{1cm}{0.4pt} &
        \rule{1cm}{0.4pt} &
        \rule{1cm}{0.4pt}
    \end{tabular}
    \end{center}

    \item Write a formula for the mass after $t$ days.

    \vspace{1cm}

    \item After how many days is only 10 g left?

    \vspace{1cm}
\end{enumerate}

\subsection*{Exercise 4: Evaluating Logarithms}

Calculate.

\begin{enumerate}
    \item $\log_2(8)=$ \rule{2cm}{0.4pt}
    \item $\log_2(32)=$ \rule{2cm}{0.4pt}
    \item $\log_{10}(1000)=$ \rule{2cm}{0.4pt}
    \item $\log_5(25)=$ \rule{2cm}{0.4pt}
    \item $\log_3(81)=$ \rule{2cm}{0.4pt}
\end{enumerate}

\subsection*{Exercise 5: Rewrite as Exponentials}

Rewrite each logarithm as an exponential equation.

\begin{enumerate}
    \item $\log_2(16)=4$
    \vspace{1cm}

    \item $\log_5(125)=3$
    \vspace{1cm}

    \item $\log_{10}(100)=2$
    \vspace{1cm}
\end{enumerate}

\subsection*{Exercise 6: Rewrite as Logarithms}

Rewrite each exponential equation as a logarithm.

\begin{enumerate}
    \item $2^6=64$
    \vspace{1cm}

    \item $3^4=81$
    \vspace{1cm}

    \item $10^5=100000$
    \vspace{1cm}
\end{enumerate}

\subsection*{Exercise 7: Solve the Equations}

\begin{enumerate}
    \item $2^x=32$

    \vspace{1cm}

    \item $10^x=1000$

    \vspace{1cm}

    \item $3^x=81$

    \vspace{1cm}

    \item $\log_2(x)=5$

    \vspace{1cm}

    \item $\log_{10}(x)=4$

    \vspace{1cm}
\end{enumerate}

\subsection*{Exercise 8: Word Problems}

\begin{enumerate}
    \item A savings account contains €500 and grows by 5\% every year.
    \begin{enumerate}
        \item Write a formula for the balance after $t$ years.
        \item How much money is in the account after 4 years?
    \end{enumerate}

    \vspace{1cm}

    \item A rumor spreads through a school. The number of students who know it triples every day. On the first day, 4 students know the rumor.
    \begin{enumerate}
        \item Write an exponential model.
        \item How many students know the rumor after 5 days?
    \end{enumerate}
\end{enumerate}

\subsection*{Challenge}

Without using a calculator, solve:

\[
2^{x+1}=64.
\]

\vspace{1cm}

\[
\log_2(x-1)=3.
\]

\vspace{1cm}

Explain in your own words why logarithms are called the \emph{inverse} of exponential functions.

\subsection*{Exercise 9: The Natural Exponential Function}

Recall that
\[
\exp(x)=e^x.
\]

Calculate the following. Give exact answers where possible.

\begin{enumerate}
    \item $\exp(0)=$ \rule{2cm}{0.4pt}
    \item $\exp(1)=$ \rule{2cm}{0.4pt}
    \item $\exp(2)=$ \rule{2cm}{0.4pt}
    \item $\exp(-1)\approx$ \rule{2cm}{0.4pt}
\end{enumerate}

\vspace{0.5cm}

Rewrite each expression using powers of $e$.

\begin{enumerate}
    \item $\exp(3)$
    \vspace{0.8cm}

    \item $\exp(-2)$
    \vspace{0.8cm}

    \item $\exp(x)$
    \vspace{0.8cm}
\end{enumerate}

\subsection*{Exercise 10: The Natural Logarithm}

Recall that
\[
\ln(x)=\log_e(x).
\]

Calculate.

\begin{enumerate}
    \item $\ln(e)=$ \rule{2cm}{0.4pt}
    \item $\ln(e^3)=$ \rule{2cm}{0.4pt}
    \item $\ln(1)=$ \rule{2cm}{0.4pt}
    \item $\ln(e^{-2})=$ \rule{2cm}{0.4pt}
\end{enumerate}

\vspace{0.5cm}

Rewrite each equation.

\begin{enumerate}
    \item $e^3=20.09\ldots$
    \vspace{1cm}

    \item $\ln(x)=2$
    \vspace{1cm}

    \item $e^x=7$
    \vspace{1cm}
\end{enumerate}

\subsection*{Exercise 11: Solving Equations}

Solve the following.

\begin{enumerate}
    \item $e^x=e^5$

    \vspace{1cm}

    \item $\ln(x)=4$

    \vspace{1cm}

    \item $e^{2x}=e^8$

    \vspace{1cm}

    \item $\ln(e^{-5})$

    \vspace{1cm}

    \item $e^{\ln(7)}$

    \vspace{1cm}

    \item $\ln(e^{12})$

    \vspace{1cm}
\end{enumerate}


\subsection*{Exercise 12}

Typical loudness levels are

\begin{center}
\begin{tabular}{l|c}
Whisper & 30 dB\\
Conversation & 60 dB\\
Busy street & 80 dB\\
Rock concert & 110 dB
\end{tabular}
\end{center}

Answer the following.

\begin{enumerate}
    \item How many times greater is the intensity of a conversation than a whisper?

    \vspace{1cm}

    \item How many times greater is the intensity of a busy street than a conversation?

    \vspace{1cm}

    \item How many times greater is the intensity of a rock concert than a whisper?

    \vspace{1cm}

    \item What is the sound pressure for a whisper?
    \vspace{1cm}
    \item What is the sound pressure for a conversation?
    \vspace{1cm}
    \item What is the sound pressure for a busy street?
    \vspace{1cm}
    \item What is the sound pressure for a rock concert?
    \vspace{1.5cm}
\end{enumerate}

\subsection*{Exercise 13}
\begin{itemize}
    \item What pressure amplitude corresponds to 0 dB? Does this mean there is no sound at all?
    \vspace{1cm}
    \item What pressure amplitude corresponds to 80 dB?
    \vspace{1cm}
    \item Assume one vacuum cleaner produces a sound with 70 dB. How loud will it be if we turn on a second vacuum cleaner? How about ten vacuum cleaners?
    \vspace{1cm}
\end{itemize}

\subsection*{Exercise 14}
Rewrite the formula for decibels with the natural logarithm instead of $\log_{10}$.

\begin{equation}\label{eq:dB}
    20 \log_{10}(\frac{p_1}{p_{\text{ref}}}) \text{dB}
\end{equation}
\vspace{4cm}

\subsection*{Exercise 15 (Challenge)}
Imagine you have a sound source that radiates sound equally in all directions. Standing 10 metres away from it, the loudness is 60 dB. How loud is it at 5 metres distance from the source? What about 20 metres?

\textit{Hint: It is easier to think about the intensity than about the pressure here. Remember: Intensity is energy per unit time per unit area. The sound source is transporting energy to its surroundings. The energy must be conserved. We can assume that the energy is distributed evenly on the surface of a sphere with a certain radius. What happens as we increase or decrease the radius of the sphere? How does does this affect the intensity? Now plug this into the dB formula below. It is the same as for the powers, since power and intensity are proportional.}

\begin{equation}
    10 \log_{10}(\frac{I_1}{I_2}) \text{dB}
\end{equation}

\end{document}